\clearpage
\thispagestyle{empty}
\pdfbookmark[0]{Observations of Magnetic Inclination in Göttingen. Edition and sources}{inclination1841-edition}
\begin{center}\Large\textbf{Observations of Magnetic Inclination in Göttingen}\end{center}

Carl Friedrich Gauss’s paper appeared in the \textit{Resultate aus den Beobachtungen des magnetischen Vereins} for 1841, article II. All 34 articles and six historical notes are AI-translated from the German text of \textit{Werke}, vol.~V (Göttingen, 1867), pp.~444--492, with AI-produced typesetting. The observations, tables and numerical notation follow that edition.

\medskip\noindent
\href{https://gdz.sub.uni-goettingen.de/id/PPN236006339}{Collected German edition (1867)}\par
\noindent\href{https://doi.org/10.1007/978-3-642-49319-5_22}{German reprint (1877)}\par
\noindent\href{https://collections.whipplemuseum.cam.ac.uk/objects/12342/}{Related Robinson dip circle, Whipple Museum, accession 3781}

\medskip\noindent
Sources cited by Gauss include Johann Caspar Horner’s \href{https://api.digitale-sammlungen.de/iiif/image/v2/bsb10131088_00105/full/1500,/0/default.jpg}{article on the dip circle}, the report by Humphrey Lloyd, Edward Sabine and James Clarke Ross on \href{https://archive.org/download/reportofbritisha05unse/page/n207.jpg}{magnetic observations in Ireland}, and Adolph Erman’s \href{https://resolver.sub.uni-goettingen.de/purl?PPN329604732}{physical observations}. The closing comparison cites \href{https://books.google.com/books?id=O_kOAAAAQAAJ&pg=PA152}{Alexandre de Humboldt’s Göttingen entries}, \href{https://doi.org/10.1017/S008045680002980X}{James D. Forbes’s measurements} and \href{https://gdz.sub.uni-goettingen.de/id/PPN35283028X_0003_2NS}{Johann Tobias Mayer’s dip-needle study}.
\clearpage
